% mage-auto-heading: foundation-model-and-agentic-harness
% mage-auto-heading: probabilistic-work-needs-external-enforcement
% Authoritative LaTeX manuscript; Markdown is a generated compatibility projection.

% <!-- part-title: The New Engineering Problem -->

% <!-- chapter-title: Agentic Machinery and Its Properties -->

\chapter{Agentic Machinery and Its Properties}

% <!-- section-terms: engineering-substrate, foundation-model, agentic-harness -->

% <!-- point: code-inspection-experiment-motivates-machinery | The DocAble experiment sought informed control over agent-produced implementation with almost no personal code inspection; this motivates examining the machinery's capabilities and limitations. | terms: engineering-substrate -->

In developing DocAble, I deliberately inspected almost none of the code the agents produced. This was partly a practical necessity, given the volume of implementation, and partly an engineering experiment: could I retain informed control over a system without personally examining its implementation? The answer depended on understanding what the agents could do, where their limitations lay, and how their work could be directed and checked.

% <!-- point: engineering-methods-fit-the-machinery | Engineering methods must fit the machinery they govern; the foundation model supplies reasoning and the harness turns that reasoning into action. | terms: engineering-substrate, foundation-model, agentic-harness -->

Engineering methods must fit the machinery they govern. Let us therefore examine the two principal components of an agent: the foundation model that supplies reasoning and the harness that turns that reasoning into action.

\section{Foundation Model and Agentic Harness}\label{heading:foundation-model-and-agentic-harness}

% <!-- section-terms: engineering-substrate, foundation-model, agentic-harness, engineered-environment -->

% <!-- point: model-plus-harness-form-a-new-engineering-substrate | A foundation model supplies inexpensive reasoning; the harness supplies runtime facilities, including hooks; the engineered environment puts those facilities to use and adds project structures beyond the harness. | terms: engineering-substrate, foundation-model, agentic-harness, engineered-environment -->

% <!-- index-def: engineering-substrate -->

A \textbf{foundation model} makes broad semantic reasoning available at a marginal cost far below the skilled human labor that many such tasks previously required. Training a model can be expensive, but using an already-trained model for one additional reasoning task is comparatively cheap.\mageaside{Engineering predates machine learning in its use of models as purposeful representations of systems. Machine learning later adopted the same word for a learned approximation of a function---also reasonably a model, but a different kind. This book gives the engineering usage precedence: unqualified model means a purposeful representation of an engineered system or some property of it. I therefore use foundation model for the trained machine-learning system and agent for the acting system built around one through a harness, tools, and working state. Where the distinction matters, I will not ask the reader to infer which ``model'' I mean.}

% <!-- index-def: agentic-harness -->

An \textbf{agentic harness} is the runtime that turns the foundation model’s reasoning into action. It manages the context supplied to the model, invokes tools in response to model outputs, and returns their results for further reasoning. It can also expose hooks: defined points at which additional code can run, for example before a tool executes.

% <!-- index-def: engineered-environment -->

The \textbf{engineered environment} is the project’s deliberate arrangement of representations, tools, processes, permissions, and checks around that machinery. It includes how the engineer configures and uses the harness’s facilities, as well as structures outside the harness. A hook facility belongs to the harness; installing a check through it to prevent edits to generated files is part of engineering the environment. Build checks and review procedures also belong to that environment even when they operate independently of the harness. The same harness can therefore support different engineered environments, depending on how engineers put its facilities to use.

% <!-- point: four-foundation-model-properties | Four properties of the foundation model matter to the engineering argument: broad semantic reasoning, probabilistic output, bounded working state, and inexpensive repeated reasoning. | terms: foundation-model -->

Four properties of the foundation model matter to the engineering argument:

\begin{itemize}
\tightlist
\item
  \textbf{Broad semantic reasoning.} A foundation model can interpret, synthesize, and search across open-ended domains. Give it an abstraction---a schema, state table, or policy---and it can reason about that abstraction rather than merely execute a fixed script.
\item
  \textbf{Probabilistic output.} A foundation model can resolve consequential engineering questions incorrectly, and its confidence in its own result is not independent evidence that the result is correct. The producing run cannot establish engineering assurance merely by affirming its own output.
\item
  \textbf{Bounded working state.} Active reasoning state is finite. Durable system knowledge must therefore persist outside the producing run and be supplied, retrieved, or reconstructed when later work needs it.
\item
  \textbf{Inexpensive repeated reasoning.} Search, generation, repair, and reconciliation---reasoning work that once consumed substantial engineering labor---can now be repeated at previously impractical scale.
\end{itemize}

% <!-- point: three-harness-capabilities | The harness contributes three complementary capabilities: action (tools turn outputs into operations), context management (assembling and maintaining working state), and interposition (points where the environment can observe, constrain, or deny an action). | terms: agentic-harness, injection-point -->

The harness contributes three complementary capabilities:

\begin{itemize}
\tightlist
\item
  \textbf{Action.} Tools turn model outputs into consequential operations: editing a file, running a command, calling an API, or changing some other part of the engineering system.
\item
  \textbf{Context management.} The harness assembles and maintains the working state available to the reasoner through mechanisms such as retrieval, memory, compaction, and task state.
\item
  \textbf{Interposition.} Because consequential actions pass through identifiable interfaces, the harness can expose points where the environment observes, constrains, redirects, or denies them.
\end{itemize}

% <!-- point: the-substrate-framing-is-an-emerging-consensus | None of these properties is alien to human engineering; what changes is their combination with broad inexpensive reasoning and a machine-mediated environment, and recent harness-engineering work treats autonomous capability as a property of the model–harness–environment system. | terms: engineering-substrate, agentic-harness -->

None of these properties is wholly alien to human software engineering. Humans also reason with bounded state, make uncertain judgments, act through tools, and rely on evidence outside themselves. What changes is their combination with broad, inexpensive semantic reasoning and a machine-mediated environment around its actions. Recent work on harness engineering similarly treats autonomous capability as a property of the foundation-model--harness--environment system rather than the foundation model alone\citep{zhong-zhu2026}\citep{young-anthropic2025}\citep{tan-thinharness2026}. \autoref{agentic-substrate} shows the resulting substrate. This combination changes the economics of familiar engineering problems and exposes new places to act on them.

\begin{figure}[htbp]
\centering
\includegraphics[width=\linewidth]{book/latex/assets/agentic-substrate-f28785289f.pdf}
\caption{\emph{Agentic machinery and environment.} The foundation model supplies reasoning, the harness turns it into tool-mediated action, and the engineered environment supplies representations, evidence, and controls around that work.}
\label{agentic-substrate}
\end{figure}

\section{The Reasoning Horizon}\label{heading:the-reasoning-horizon}

% <!-- section-terms: long-horizon-task, reasoning-horizon, churn, context-window -->

% <!-- point: scale-creates-the-reasoning-problem-before-an-agent-enters | Human working memory is limited even when long-term knowledge is extensive; notes, diagrams, and other external representations let engineers work with more information than they can hold actively in mind. | terms: long-horizon-task, context-window -->

Scale has always created problems. An engineer may be able to hold a small program in mind, but a large system soon exceeds that capacity. Human \textbf{working memory} is limited: we can actively hold and work with only a small amount of information at once, even when we have much more knowledge in long-term memory\citep{cowan2001capacity}. A change that looks simple in one file may require us to recall an interface elsewhere, a prior design decision, and a constraint on acceptable behavior. We use notes, diagrams, and other external representations so that we do not have to keep all of that information in our heads\citep{norman2013}.

% <!-- point: once-a-task-exceeds-active-state-question-relevant-state-must-be-reconstructed | A harness can preserve and retrieve information across invocations, extending work beyond the model's immediate context; retaining a record does not ensure that the agent will find and use it. | terms: reasoning-horizon, churn -->

% <!-- index-def: long-horizon-task -->

An agent encounters a comparable limit through its foundation model's context window: the bounded input available in a single invocation. The harness may store a plan, earlier decisions, or tool results outside that window and retrieve them later. That lets work continue across many invocations, but information needed for a particular judgment must be supplied when that judgment is made. A \emph{long-horizon task} requires a sequence of related decisions, with earlier results or commitments carried into later steps. Merely retaining a record does not ensure that the agent will find and use it.

Memory management can nevertheless expand what the whole system can work over. Research on agent memory and context management has demonstrated this through external storage, summarization, and selective retrieval\citep{packer2023memgpt}\citep{lee2024readagent}. The model's context window stays bounded, but the surrounding machinery makes more information usable across the task.

% <!-- point: availability-is-not-a-useful-representation | Available information is not always used effectively: usable context can fall short of nominal capacity, information within a window is not uniformly accessible, and repository work incurs navigation and reconstruction costs. | terms: reasoning-horizon, context-window -->

Having information available does not guarantee that it can be used effectively. Studies of long-context language models find that usable context can fall short of nominal capacity and that information within the context is not uniformly accessible\citep{hsieh2024ruler}\citep{liu2024lostmiddle}. Repository-scale studies likewise report substantial navigation and reconstruction work as relevant state spreads across a codebase\citep{sourcegraph2026largecodebase}.

% <!-- point: the-reasoning-horizon-defined | The reasoning horizon is the practical limit of work the reasoner and its supporting machinery can carry through reliably, given the available representations and resources; task demands are distinct from this system-level limit. | terms: reasoning-horizon, context-window -->

% <!-- index-def: reasoning-horizon -->

In this book, we use the term \emph{reasoning horizon} to describe the practical limit of the work a reasoner and its supporting machinery can carry through reliably, given the available representations and resources. For an agent, this depends on the foundation model together with the harness, external memory, and tools.\footnote{We borrow the term from Ramezanali and colleagues, who use it in seqBench to describe the length of a sequence a model can reliably complete\citep{ramezanali2025seqbench}. Here we extend it to the practical limits of the complete reasoning system, including its harness, external memory, and tools.}

The reasoning horizon can expand in two ways: better machinery can preserve more relevant state, while better representations can reduce how much state the task requires. Research on external representations has long shown that how information is organized affects the work needed to reason with it\citep{larkin1987diagram}\citep{zhang1994representations}. These improvements may also make work cheaper without extending the range of tasks the system can handle. \autoref{reasoning-horizon-fig} illustrates how purposeful abstraction reduces reconstruction.

\begin{figure}[htbp]
\centering
\includegraphics[width=\linewidth]{book/latex/assets/reasoning-horizon-7eba103f30.pdf}
\caption{\emph{Representation and the reasoning horizon.} When relevant state cannot all remain active, it must be preserved elsewhere and recovered. Raw implementation can demand repeated reconstruction; purposeful abstraction can reduce that burden, making work cheaper and sometimes bringing more demanding tasks within the system's reasoning horizon.}
\label{reasoning-horizon-fig}
\end{figure}

% <!-- box-family: inset -->

% <!-- inset-domain: AGENT HARNESS -->

% mage-box-prefix: Inset —

\begin{bookbox}{inset}{Context Windows and Compaction}
\label{inset-context-windows-and-compaction}
Intuitively, a harness can extend an agent's reasoning horizon by preserving useful information and making it available when needed. Context windows and compaction make that idea concrete for agents built around large language models (LLMs). A \textbf{context window} limits how much material a model can process in one invocation. This capacity is measured in tokens. For text, tokens may be words, parts of words, or punctuation; multimodal models also encode images, audio, and other inputs into representations that occupy space in the context. Instructions, earlier exchanges, selected file contents, and tool results occupy space in the window, with room also needed for the response. A long-running task can accumulate far more history than one window will hold.

The harness can make room by compacting that history. \textbf{Compaction} replaces part of the accumulated context with a shorter account, often produced by asking an LLM to summarize it. Later invocations receive that account alongside selected original material and new input. For example, pages of investigation and test output might become a handoff recording the current plan, decisions already made, checks completed, and work still open.

Compaction can lose distinctions that later matter. If a summary preserves ``the tests passed'' but drops which tests ran or which version they exercised, the next invocation may treat that result as stronger evidence than it is. Records kept outside the window can supply the missing detail, provided the harness or agent knows to retrieve them. The engineering problem is choosing what must survive compaction and what can safely be recovered later\citep{rajasekaran2025contextengineering}.

The generation of commodity intelligence that motivated this book is built around language models with bounded context windows. We deliberately abstract from that machinery: the \textbf{reasoning horizon} concerns what a reasoner and its supporting machinery can carry through reliably, given the available representations and resources. The concept is thus broader than a context-window limit, and is intended to remain useful as new kinds of intelligence bring different capabilities and constraints.
\end{bookbox}

\section{Probabilistic Work Needs External Enforcement}\label{heading:probabilistic-work-needs-external-enforcement}

% <!-- section-terms: foundation-model, engineering-substrate -->

% <!-- point: authorship-is-separate-from-assurance | Software engineering has always separated authorship from assurance — we compile instead of trusting a prediction, test instead of accepting an explanation — because the producer's confidence is not enough. | terms: engineering-substrate, foundation-model -->

The second problem concerns assurance. Software engineering has always separated authorship from assurance. We compile code to check that it compiles and run tests to evaluate whether it behaves as required. Consequential systems add review, static analysis, formal methods, operational evidence, or independent approval. The producer\textquotesingle s confidence is insufficient evidence of correctness.

% <!-- point: enforcement-must-reside-outside-the-agent | Commodity intelligence makes that separation more important because implementation is cheap and repeatable; an agent can propose, inspect, explain, and repair, but another judgment from the same substrate does not constitute independent evidence — the environment must establish what evidence counts. | terms: engineering-substrate, foundation-model -->

Cheap, fast, repeatable agentic production makes that separation more consequential. An agent can propose a change, inspect it, explain it, and often repair it, but another probabilistic judgment from the same productive substrate does not automatically create engineering assurance. If an obligation matters, the environment must establish what evidence counts and what follows when the obligation is not met.

The U.S. defense community\textquotesingle s AI4SDLC guidance identifies the same scaling problem: as agents produce more work asynchronously, the people reviewing it can become the bottleneck\citep{codeMilWorkflowGovernance2026}. Requiring human approval is therefore only part of the answer; reviewers also need representations and evidence that let them make an informed judgment about what the machinery has produced.

Commodity intelligence is probabilistic. This book focuses on the implications for engineering. A foundation model's answer does not ordinarily come with a guarantee that its engineering judgments are correct. There are several sources of uncertainty: generation can vary between runs, a request may admit several interpretations, and the model may take different reasoning steps or choose different tools. A task can depend on many such judgments before it produces a result. The reliability of this work also depends on how the agent's harness and engineered environment organize it. They can arrange retries, decompose the task, combine multiple answers, retrieve additional context, or ask another model to evaluate the first. These techniques can improve reliability, but none by itself turns an inferred judgment into an engineering guarantee.

A simple model shows how an agent's reliability on individual features affects its chance of implementing them all correctly. Suppose an agent implements \(n\) features in a product, numbered 1 through \(n\). Let \(p_i\) denote the probability that the agent implements feature \(i\) correctly. Under the simplifying assumption that these implementation outcomes are independent, the probability that the agent implements all the features correctly is

\[
P(\text{agent succeeds on all features}) = \prod_{i=1}^{n} p_i
\]

Remember, when you multiply fractions between zero and one, the product is smaller than any of the factors. The effect is small when the agent is highly reliable and has only a few features to implement. If an agent has a high probability of correctly implementing each feature, say 99\%, then its probability of getting three features right is \(0.99^3 \approx 97\%\). As we ask the agent to implement more features, its chance of getting them all right falls: \(0.99^{10} \approx 90\%\) for ten features, and \(0.99^{100} \approx 37\%\) for a hundred. If the agent is less reliable, say 95\% on each feature, its chance of correctly implementing all ten falls to \(0.95^{10} \approx 60\%\). An agent that usually gets individual features right can therefore have a much lower chance of getting every feature in a product right.

% <!-- index-def: probabilistic-surface -->

% <!-- gloss-only: Probabilistic surface | The decisions in a system that depend on agentic inference. A larger probabilistic surface creates more opportunities for inference error to affect the outcome, even when each individual decision is highly reliable. MAGE seeks to reduce this surface by moving suitable decisions into explicit models, deterministic mechanisms, and independently enforced obligations, while leaving genuinely judgment-dependent decisions probabilistic. -->

% <!-- point: the-probabilistic-surface-is-the-judgments-still-bought-per-instance | Decisions that depend on agentic inference form the probabilistic surface; mechanisms can enforce particular requirements even when the agent makes a mistake. | terms: engineering-substrate, foundation-model -->

The decisions that depend on agentic inference form the system's \textbf{probabilistic surface}. Other parts of the system may be deterministic, such as a tool that applies a fixed transformation or a test suite that evaluates explicit conditions. Better models and iterative loops can make the agent's judgments more reliable. As we will discuss throughout the book, we can also build mechanisms that enforce particular requirements even when the agent makes a mistake.

On the other hand, even if agentic implementation were perfectly reliable, that would not tell us what to build. Engineers must still choose and justify the requirements, and determine what validation evidence would establish that the resulting system meets the needs it is meant to serve.
